% Note: several 2025--2026 citation keys below are inferred from the local
% review notes and may need reconciliation with the thesis bibliography file.
\chapter{Literature Review}\label{chap:survey}

This chapter reviews the literature most relevant to efficient and robust GIS-scale 3D reconstruction from remote sensing imagery. It first presents a systematic mapping of the local corpus, spanning foundational photogrammetry, satellite stereo and digital elevation model (DEM) production, recent deep learning approaches for acceleration, and emerging work on modular and adaptive reconstruction workflows. It then synthesizes the literature with respect to the three research questions of this thesis: acceleration through deep learning, robustness through modularity, and adaptive workflow management. The central argument emerging from the corpus is that computational efficiency and generalization should be treated as coupled pipeline design problems rather than as isolated model-performance problems.

\section{Systematic Literature Review and Mapping}

\subsection{Scope, objectives, and review strategy}

The review focuses on image-based 3D reconstruction methods intended to produce GIS-usable outputs such as georeferenced point clouds, surface models, and DEMs from large remote sensing image collections. The scope covers structure-from-motion (SfM), multi-view stereo (MVS), bundle adjustment (BA), satellite stereo pipelines, neural surface and neural field methods, and workflow-level strategies for scaling reconstruction across heterogeneous scenes. Three objectives guided the review: (i) to identify where end-to-end runtime is spent, (ii) to determine how recent work exposes additional parallelism or reduces optimization cost, and (iii) to assess how robustness is maintained when scene conditions, sensors, or acquisition times vary.

The review strategy followed two passes over the project-curated local corpus. The first pass retained foundational studies when they established baselines that still shape current reconstruction pipelines or revealed bottlenecks that remain relevant at GIS scale. The second pass retained recent studies when they contributed at least one of the following: explicit acceleration of a computationally expensive reconstruction stage, a modular or hybrid architecture designed for heterogeneous conditions, or an adaptive selection mechanism that is relevant to autonomous workflow control. Studies were deprioritized when they focused on unrelated visual tasks, lacked methodological detail, or did not provide transferable insight for remote sensing reconstruction.

\subsection{Thematic mapping of the corpus}

Table~\ref{tab:litmap-summary} summarizes the main thematic clusters identified in the reviewed literature.

\begin{table}[ht]
\centering
\caption{Systematic mapping of representative literature on GIS-scale 3D reconstruction.}
\small
\begin{tabular}{p{0.19\textwidth}p{0.27\textwidth}p{0.16\textwidth}p{0.25\textwidth}}
\hline
\textbf{Cluster} & \textbf{Representative studies} & \textbf{Primary unit of computation} & \textbf{Main signal for this thesis} \\
\hline
Classical photogrammetry and satellite stereo baselines & \cite{Schonberger2016,Schonberger2016MVS,Galliani2015,deFranchis2014} & Tracks, pixels, stereo tiles & Establishes high-quality baselines and shows that dense matching and BA remain persistent bottlenecks \\
GPU-first and feed-forward acceleration & \cite{Li2025FastMap,Zhong2025InstantSfM,Yu2025CuSfM,Safari2025,Gopinath2025,Light3RSfM2025,DenseSfM2025} & Image pairs, sparse blocks, dense correspondences & Runtime gains come from reformulating SfM and BA around GPU-friendly sparse or feed-forward primitives \\
Scalable neural surface and field reconstruction & \cite{SatDN2025,ShadowGS2026,SnakeNeRF2025,BirdNeRF2025} & Rays, tiles, local fields, hash-grid blocks & Learned surface models can reduce per-scene optimization cost, but scalability depends on tiling, priors, and stable training behavior \\
Robust and modular remote sensing workflows & \cite{Amadei2025,Masquil2026,FoundationStereo2025,MVSR3D2025,Planet4Stereo2025} & Pipeline stages, branches, domain-specific modules & Robustness is improved by combining classical and learned stages instead of expecting one model to generalize across all conditions \\
Adaptive supervision and supervisory data & \cite{SMoEStereo2025,MoE3D2026,SatDepth2025,SEO2025,Elias2025} & Experts, scene diagnostics, metadata-conditioned policies & Explicit routing, diagnostics, and quality-gated adaptation are emerging, but workflow-level supervisors remain underdeveloped \\
\hline
\end{tabular}
\label{tab:litmap-summary}
\end{table}

The mapping suggests four broad trends. First, classical photogrammetric pipelines remain the reference point for geometric reliability and GIS-grade outputs, especially where georeferencing and DEM validity matter~\cite{Schonberger2016,Schonberger2016MVS,deFranchis2014}. Second, many strong acceleration results are achieved not by simply moving existing code onto GPUs, but by changing the computational structure itself, for example through pair-wise formulations, sparse block scheduling, mixed-precision BA, or feed-forward alignment~\cite{Li2025FastMap,Zhong2025InstantSfM,Safari2025,Light3RSfM2025}. Third, remote-sensing-specific deep learning increasingly depends on explicit handling of time, shadows, sensor geometry, and scale rather than assuming that methods developed for ground imagery will transfer directly~\cite{Masquil2026,SatDepth2025,SEO2025,ShadowGS2026}. Fourth, the literature is shifting from monolithic model thinking toward system-level design, where modular decomposition, routing, priors, and quality gates together determine whether a workflow is both efficient and operationally robust~\cite{MVSR3D2025,SMoEStereo2025,MoE3D2026,Elias2025}.

\section{Literature Synthesis in Relation to the Research Questions}

\subsection{RQ1: Acceleration through deep learning}

\textbf{RQ1.} How can deep learning formulations be integrated into remote sensing reconstruction pipelines to significantly reduce end-to-end computational runtime for GIS-scale datasets?

The literature is consistent that runtime bottlenecks are concentrated in a limited set of stages: image matching, dense depth estimation, BA, and repeated per-scene optimization for surface inference~\cite{Schonberger2016,Schonberger2016MVS,Galliani2015}. This is important because it suggests that full pipeline replacement is not required to obtain substantial runtime gains. Instead, many recent contributions accelerate reconstruction by replacing or restructuring the most expensive stages while leaving the broader geometric workflow intact.

Within sparse reconstruction, recent work shows a shift from iterative global optimization toward feed-forward or GPU-native formulations. FastMap, InstantSfM, and CuSfM all expose more concurrency by reorganizing SfM and BA into units that are better aligned with accelerator hardware, such as pair-structured computations, sparse blocks, or graph-based factors~\cite{Li2025FastMap,Zhong2025InstantSfM,Yu2025CuSfM}. Matrix-free and mixed-precision BA work further shows that solver design is itself a major source of runtime reduction, especially when memory pressure is the real bottleneck rather than raw floating-point throughput~\cite{Safari2025,Gopinath2025}. Feed-forward approaches extend this logic: Light3R-SfM reduces dependence on heavyweight iterative alignment through learnable global alignment and retrieval-guided graph construction, while Dense-SfM pushes more of the matching burden into dense correspondences that can improve track quality before downstream optimization~\cite{Light3RSfM2025,DenseSfM2025}.

Within dense reconstruction, the literature supports deep learning as a practical acceleration mechanism when dense estimation is naturally parallelizable on GPUs. Satellite stereo papers based on transformer-CNN feature fusion, self-supervised satellite-adapted features, and RAFT-like disparity estimation show that dense correspondence estimation can benefit substantially from learned representations, especially in weak-texture or repetitive-image regions where traditional matching struggles~\cite{Amadei2025,FoundationStereo2025,SatDepth2025}. However, the evidence also shows that acceleration is not only about raw inference speed. It also depends on whether the learned stage reduces downstream failure and reprocessing. In other words, a faster depth network that produces unstable geometry under temporal or sensor variation may not improve end-to-end throughput in practice.

Neural field and neural rendering approaches reinforce this point. Methods such as Sat-DN, ShadowGS, Snake-NeRF, and BirdNeRF are motivated in part by the high optimization cost of per-scene neural reconstruction, and they reduce this cost through progressive training, multi-resolution encodings, shadow-aware constraints, and explicit tiling or out-of-core training~\cite{SatDN2025,ShadowGS2026,SnakeNeRF2025,BirdNeRF2025}. Yet these methods do not eliminate systems concerns; instead, they shift the acceleration problem toward scene partitioning, prior injection, and training stability at large geographic extents. The literature therefore supports a selective replacement strategy for RQ1: deep learning is most effective when it is inserted at computational choke points and tightly integrated with pipeline-level scheduling and quality control, rather than treated as a complete substitute for all classical geometry.

\subsection{RQ2: Robustness through modularity}

\textbf{RQ2.} To what extent can a modular reconstruction architecture improve robustness and adaptability compared to monolithic learning-based approaches?

The strongest motivation for modularity is the repeated observation that generalization failures in remote sensing are structured rather than random. Performance degrades for identifiable reasons: changes in sensor geometry, revisit time, illumination, land cover, shadowing, or scene composition~\cite{Masquil2026,SatDepth2025,SEO2025}. This matters because a monolithic model trained to absorb all such variation implicitly must solve multiple domain shifts at once, often with little interpretability about which conditions caused failure. The recent literature instead trends toward modularity, whether explicitly at the workflow level or implicitly within network design.

Several lines of evidence support a modular architecture over a single universal model. First, operational photogrammetric pipelines still rely on stage decomposition because it enables targeted replacement while preserving geometric control. Classical SfM, MVS, and satellite stereo remain competitive baselines precisely because individual stages can be improved independently without destabilizing the entire workflow~\cite{Schonberger2016,Schonberger2016MVS,deFranchis2014,Planet4Stereo2025}. Second, hybrid systems increasingly combine learned and classical components rather than choosing one paradigm exclusively. This is visible in workflows that use learned matching, depth priors, or semantic cues while retaining explicit geometry, triangulation, and map-oriented merging~\cite{SatDN2025,MVSR3D2025,Elias2025}. Third, large-scale scaling strategies themselves are modular: tile-based processing, local field decomposition, and clustered sub-scenes turn large-area reconstruction into a collection of bounded subproblems that can be solved, checked, and merged more reliably~\cite{SnakeNeRF2025,BirdNeRF2025,Zhou2026}.

The literature also suggests that modularity improves robustness because it supports algorithmic diversity. Under this view, the same pipeline stage may have multiple candidate implementations: a classical stereo matcher for well-behaved same-date imagery, a learned stereo model for weak texture, a diachronic method when appearance drift is severe, or a neural field method when view sparsity and shadowing dominate~\cite{Masquil2026,ShadowGS2026,FoundationStereo2025}. Modular architecture therefore acts as a mechanism for conditional specialization. It allows each stage to use the method most appropriate to the current data regime instead of forcing one model to remain optimal everywhere.

Importantly, the literature does not imply that modularity is free. Decomposition introduces interfaces, merge logic, and consistency risks. Distributed BA and hierarchical reconstruction show that local robustness does not automatically produce global correctness; sub-model alignment and merge consistency become critical failure surfaces at GIS scale~\cite{Zheng2023,Li2025BALM,Zhou2026}. Nevertheless, this is still a strong argument for modular design, because those interfaces can be measured, validated, and improved explicitly. By contrast, failure inside a monolithic model is often harder to diagnose or localize. Overall, the literature supports RQ2 by indicating that modular architectures offer a better path to robust generalization, provided that stage boundaries and merge contracts are designed deliberately.

\subsection{RQ3: Adaptive workflow management}

\textbf{RQ3.} How can an adaptive supervisory mechanism effectively select and configure reconstruction modules in response to varying scene characteristics and data conditions?

Compared with acceleration and modularity, the literature on explicit workflow supervision is less mature, but the underlying ingredients are increasingly visible. The first ingredient is conditional expert selection. Selective mixture-of-experts methods such as SMoEStereo and MoE3D show that routing can be learned at scene level or even pixel level so that different structures are processed by different specialists~\cite{SMoEStereo2025,MoE3D2026}. Although these papers embed supervision inside the model, they are conceptually aligned with a pipeline-level supervisor: both treat heterogeneity as a selection problem rather than as noise to be averaged away.

The second ingredient is diagnostic metadata about failure conditions. Diachronic stereo work makes time gap and appearance drift explicit control variables rather than nuisance factors, while S-EO and shadow-aware reconstruction work show that shadow inconsistency can be detected and exploited to choose more appropriate constraints or reconstruction families~\cite{Masquil2026,SEO2025,ShadowGS2026}. SatDepth and related satellite-specific supervision datasets provide a further foundation for adaptive control because they make known satellite image failure modes measurable during training and evaluation~\cite{SatDepth2025}. Together, these studies suggest that a supervisory layer can use observable cues such as revisit interval, expected seasonal change, texture weakness, shadow severity, or confidence diagnostics to decide whether to deploy classical stereo, learned stereo, extra priors, fine-tuning, or an alternative reconstruction branch.

The third ingredient is workflow-level quality control. Large-scale and distributed reconstruction papers repeatedly show that failure recovery, merge validation, and consistency checks are essential once computation is partitioned~\cite{Zheng2023,Li2025BALM,Zhou2026,Elias2025}. This means a useful supervisor cannot only select modules at the start of the pipeline; it must also monitor intermediate outputs and intervene when quality deteriorates. In practice, this may include triggering fallback modules, tightening thresholds, revisiting dense matching, or rejecting locally acceptable but globally inconsistent merges.

The key gap is that most current literature still treats these supervisory ideas implicitly. Routing is usually embedded within a single model, adaptation is often paper-specific, and quality gates are reported as engineering details rather than as first-class experimental variables~\cite{SMoEStereo2025,SatDN2025,Elias2025}. Few studies evaluate an explicit, modular supervisor that jointly reasons about acceleration, robustness, and GIS-grade output quality across the full reconstruction workflow. This gap directly motivates the research direction of the thesis.

\section{Research Gaps and Thesis Positioning}

The synthesis above reveals four research gaps that shape the thesis contribution.

First, many reported speedups are stage-local rather than end-to-end. Papers often show acceleration for dense matching, SfM, BA, or neural field training individually, but much less evidence links those gains to full GIS-scale reconstruction workflows with consistent map-ready outputs~\cite{Li2025FastMap,SatDN2025,Elias2025}. This leaves open the question of how much runtime is truly saved after data preparation, fallback handling, georeferencing, and product generation are included.

Second, recent deep learning work still struggles to demonstrate stable transfer across heterogeneous geographic environments. The literature repeatedly identifies multi-date imagery, illumination change, sensor heterogeneity, sparse views, and weak texture as central sources of failure~\cite{Masquil2026,SatDepth2025,SEO2025,ShadowGS2026}. This suggests that robustness is unlikely to come from a single monolithic network and instead depends on modular specialization and controlled interaction between modules.

Third, the literature contains many ingredients for adaptive supervision but relatively few complete supervisory frameworks. Expert routing, domain adaptation, confidence-aware refinement, and shadow- or time-aware policies are all present in partial form~\cite{SMoEStereo2025,MoE3D2026,FoundationStereo2025,Masquil2026}. What is missing is a workflow-level design that turns those ideas into an explicit decision-making layer for selecting and configuring modules across the whole pipeline.

Fourth, GIS-grade evaluation criteria are still weakly integrated into the design of learned reconstruction systems. Reconstruction papers often report geometric or rendering metrics, while operational remote sensing workflows additionally require positional consistency, merge validity, reproducibility, and downstream suitability for GIS analysis~\cite{deFranchis2014,Elias2025,Zhou2026}. This creates a gap between model-centric performance reporting and operational utility.

Taken together, these gaps justify the thesis direction proposed in this research. The literature supports developing a modular 3D reconstruction framework in which deep learning is used selectively for high-cost stages, classical geometry is retained where it provides stability or interpretability, and an adaptive supervisory mechanism configures the workflow according to scene and sensor conditions. Such a design directly addresses the unresolved connection between acceleration, robustness, and operational autonomy in GIS-scale remote sensing reconstruction.

\section{Summary}

This chapter reviewed the literature on efficient and robust GIS-scale 3D reconstruction from remote sensing imagery through a systematic mapping and a research-question-driven synthesis. The evidence shows that the main computational bottlenecks remain concentrated in matching, dense estimation, BA, and repeated scene optimization; that recent deep learning advances are most effective when they selectively replace those bottlenecks rather than attempt universal end-to-end substitution; that robustness under geographic and temporal heterogeneity is better supported by modular and hybrid workflows than by monolithic models; and that the ingredients for adaptive workflow supervision are emerging but not yet integrated into complete pipeline-scale systems. These findings motivate a thesis framework centered on selective deep learning acceleration, modular pipeline design, and explicit adaptive supervision to produce fast, robust, and GIS-usable 3D reconstruction outputs.
